% AAAI-27 Student Abstract and Poster Program draft
% Track deadline: 28 September 2026
% Requirement: maximum 2 pages INCLUDING references, AAAI two-column camera-ready style.
% IMPORTANT before submission:
%   1. Compile with the official AAAI-27 author kit (aaai27.sty).
%   2. Replace the phone-number placeholder below.
%   3. Check the final PDF is no more than 2 pages including references.
%   4. Submit a separate supplemental PDF/ZIP (<=20 MB), as requested by the CFP.


\documentclass[letterpaper]{article}

% Use the official AAAI-27 Author Kit files: aaai2027.sty and aaai2027.bst.
\usepackage{aaai2027}

\usepackage[hyphens]{url}
\usepackage{graphicx}
\usepackage{booktabs}
\usepackage{microtype}
\usepackage{natbib} % Required by the AAAI 2027 bibliography style.

\urlstyle{rm}
\def\UrlFont{\rm}

% AAAI papers normally use unnumbered sections.
\setcounter{secnumdepth}{0}

\title{Evaluating Visual Evidence Dependence in Medical Vision--Language Models}

\author{
    Dawud Izza\textsuperscript{\rm 1},
    Frank Soboczenski\textsuperscript{\rm 1}
}
\affiliations{
    \textsuperscript{\rm 1}University of York, United Kingdom\\
    dawud.izza@york.ac.uk, frank.soboczenski@york.ac.uk
    % Optional: add a personal/research URL for Dawud if desired.
}

\begin{document}

\maketitle

\begin{abstract}
Medical vision--language models (VLMs) can answer questions about clinical images, but answer accuracy alone does not establish that the corresponding image was used. We present a preliminary, development-only study of visual dependence, confidence sensitivity, and attribution-intervention sensitivity in gastrointestinal (GI) endoscopy visual question answering (VQA). We fine-tune paired-image and constant-image adapters, then evaluate controlled image interventions on a locked development selection. Across two seeds, the model is sensitive to image identity, but a neutral image unexpectedly outperforms the corresponding image. This result demonstrates both the value of behavioural visual controls and the danger of interpreting confidence or plausible answers as evidence of faithful grounding. The study provides a reproducible, test-sealed protocol for distinguishing intervention sensitivity from faithful visual grounding.
\end{abstract}

\section{Motivation and Research Gap}

GI endoscopy is a visually demanding clinical setting in which relevant findings can vary substantially with anatomy, viewpoint, illumination, image quality, bowel preparation, and procedural artefacts~\cite{beg2017quality}. Recent datasets and shared tasks such as Kvasir-VQA, Kvasir-VQA-x1, and Medico provide image--question--answer pairs for evaluating multimodal reasoning over gastrointestinal imagery~\cite{gautam2024kvasirvqa,gautam2025kvasirvqa,gautam2025medico}. Parameter-efficient adaptation has also shown promise for GI VQA~\cite{peter2026parameterefficient}.

\section{Related Work and Gap}

Answer quality alone is insufficient to establish that a VLM is using the image in the way its explanation implies. Medical VLM evaluations have identified hallucination and weak visual grounding as continuing concerns~\cite{chang2025medheval,liu2026medical}. A model may exploit linguistic or dataset regularities, continue answering when visual evidence is absent, or provide an apparently coherent explanation whose highlighted evidence is not important to the prediction. In a clinical setting, such explanation--evidence mismatch is especially important because a plausible explanation may encourage reliance without demonstrating visual faithfulness.

This work therefore asks three linked questions:
\begin{enumerate}
    \item \textbf{Visual dependence:} Does the answer depend on the supplied image?
    \item \textbf{Intervention sensitivity:} Do attribution-selected regions influence the answer when perturbed?
    \item \textbf{Confidence sensitivity:} Does model confidence respond when relevant visual evidence is reduced or altered?
\end{enumerate}

\section{Evaluation Approach}

The study uses GI VQA as a controlled case study. We adapt PaliGemma using 256 training source images and evaluate four questions per source on a locked 256-source development selection, with no access to the test partition. The main comparison holds the paired-image adapter fixed while replacing the corresponding image with a shuffled image, a neutral image, a blurred version of the correct image, or a constant-gray image. A constant-image adapter provides a descriptive question-prior control, while the unadapted base model provides a reference point. Normalized token-F1 is the primary answer metric; sequence confidence, calibration, and paired confidence contrasts are secondary behavioural measures.

Attribution intervention sensitivity is evaluated separately from answer correctness. Regions selected by Grad-CAM or decoder attention are blurred or replaced, and changes in the answer or confidence are measured. These interventions test whether selected regions influence behaviour; they do not by themselves establish clinically faithful explanations.

Finally, confidence is treated as an additional behavioural signal. Prior work has shown that predictive performance and uncertainty need not be well aligned in VLMs~\cite{kostumov2024uncertainty}. The study therefore examines whether confidence decreases when relevant visual evidence is weakened, removed, or replaced, and whether this sensitivity is associated with changes in answer correctness and explanation behaviour.

\section{Current Status and Intended Contribution}

Current development focuses on a GI VQA baseline using Kvasir-VQA-x1. Two independent seeds completed the controlled evaluation. The pooled normalized token-F1 scores were 0.3605 for correct, 0.3427 for blurred-correct, and 0.3916 for both neutral and constant-gray images. Thus, blurring produced a modest reduction (correct minus blurred: $+0.0178$), while neutral exceeded correct (correct minus neutral: $-0.0311$). Neutral and constant-gray images had different hashes for every item, yet produced identical predictions and confidence values item by item.

\begin{table}[t]
\centering
\small
\caption{Pooled development performance. Fine-tuning improves over the unadapted base, but the gain is not uniquely visual.}
\label{tab:performance}
\begin{tabular}{lr}
\toprule
\textbf{Condition} & \textbf{Token-F1} \\
\midrule
Unadapted base + correct image & 0.1650 \\
Constant-image adapter + neutral & 0.3309 \\
Paired adapter + correct image & 0.3605 \\
Paired adapter + neutral image & \textbf{0.3916} \\
\bottomrule
\end{tabular}
\end{table}

The paired adapter therefore improves over the unadapted PaliGemma baseline by $+0.1955$ F1, while the constant-image adapter improves by $+0.1659$. This establishes a VQA adaptation benefit, but not a reliable benefit from the corresponding image.

\begin{table}[t]
\centering
\small
\caption{Exploratory category breakdown of correct-minus-neutral token-F1. Categories are multi-label and are not independent tests.}
\label{tab:categories}
\begin{tabular}{lrr}
\toprule
\textbf{Question category} & \textbf{Items} & \textbf{$\Delta$F1} \\
\midrule
Finding presence & 176 & $-0.3866$ \\
Instrument count & 42 & $+0.2178$ \\
Abnormality presence & 32 & $+0.2808$ \\
\bottomrule
\end{tabular}
\end{table}

Confidence was not aligned with answer performance: correct-image confidence exceeded neutral and constant-gray confidence by 0.0411 despite their higher F1, and exceeded blurred-correct confidence by 0.0105. Grad-CAM intervention effects averaged $+0.0285$ for blur deletion and $+0.0622$ for blur insertion, whereas attention effects ranged from $-0.0424$ to $+0.0276$. The largest category-specific reversal was finding presence; this exploratory result motivates further controls rather than establishing a general category effect.

As a secondary semantic check, a blinded 0/1/2 rubric was applied by GPT-4.1 to all 2,048 item-conditions and by Gemini 3.8 Flash to a deterministic 50\% subset. Their correct-minus-neutral mean label contrasts were $-0.348$ and $-0.242$, respectively, with 89.5--94.5\% exact agreement on the shared subset. The judges therefore reproduce the neutral advantage, but remain post-hoc evaluators rather than clinical adjudicators.

The paired-versus-constant comparison changes both the trained adapter and the image condition, so it is not a clean causal comparison. The shuffled, neutral, blurred, and gray conditions are the stronger same-adapter visual-ablation tests. All reported results are preliminary and development-only; the sealed test partition was not accessed. The intended contribution is a reproducible methodology for distinguishing visual sensitivity and intervention effects from faithful visual grounding in medical VLMs.

\bibliography{references}

\end{document}
